\section{Related Literature}
\label{sec:related_literature}

\subsection{Commodity prices, exchange rates and connectedness in Ghana}

The relationship between commodity markets and the Ghanaian cedi has been studied from several perspectives, but the literature has developed mostly through monthly-frequency models of asymmetric response, volatility transmission and time-frequency comovement. \citet{zankawah2020}, using monthly data from 1991 to 2015 and multivariate GARCH specifications, report significant spillover effects from crude oil prices to Ghana's exchange rate. Their analysis is important because it treats exchange-rate volatility as part of a broader commodity-shock system, but it does not examine joint tail dependence or the sensitivity of daily dependence estimates to market-calendar construction.

A later group of studies broadened the commodity set and placed greater emphasis on nonlinear relationships. \citet{archer2022} analyse cocoa, gold, Brent crude oil and the nominal exchange rate with quantile regression over September 2007 to December 2020. Their results vary across exchange-rate quantiles, which is consistent with the view that commodity-currency relationships can depend on market conditions rather than follow a single linear pattern. \citet{boateng2022} similarly document time-frequency variation in the relationships among cocoa, Brent crude oil, gold and Ghanaian macroeconomic indicators using wavelet methods. Although their focus is not the Cedi alone, their findings reinforce the broader point that commodity linkages in Ghana vary by horizon and state.

\citet{barson2023} move closer to the present question by studying commodity-return comovement in Ghana and the role of the exchange rate. Using monthly observations from September 2007 to March 2021, they find stronger coherence at short and medium horizons and conclude that the exchange rate plays an important role in commodity interdependence. \citet{yeboah2025} provide more recent evidence using quantile-on-quantile estimation with monthly data from 2003 to 2023. Their results also point to asymmetric and nonlinear relationships between commodity prices and the exchange rate. Taken together, these studies make a constant, symmetric commodity-Cedi relationship difficult to defend.

Recent work has expanded the Ghanaian literature beyond direct commodity-price regressions. \citet{atipaga2025} use daily Cedi data from June 2016 to October 2023 with wavelet and Rényi transfer entropy methods to examine global uncertainty channels. Their results associate Cedi dynamics with oil-price volatility and Ghana's sovereign spread, with effects varying across horizons. \citet{woode2025} examine time- and frequency-varying connectedness among commodities, currencies and equity markets in commodity-dependent sub-Saharan African economies, including Ghana, and report changing transmission patterns across crises and horizons. These studies confirm that frequency, market state and model choice matter for how cross-market dependence is described.

Copula methods themselves are not new to Ghana. \citet{mensah2020} use daily GHS exchange rates against the U.S. dollar, euro and pound sterling from 2009 to 2019 and estimate both static and time-varying copulas. They report positive dependence and evidence that tail dependence evolves over time. This study is an important novelty guardrail for the present paper. The contribution here is not the introduction of copulas to Ghanaian financial data. Rather, the question is whether commodity-Cedi dependence, copula adequacy and stress-period tail conclusions remain defensible under alternative treatments of nonsynchronous market calendars, explicit marginal diagnostics and multiplicity-controlled inference.

The Ghana studies reviewed here therefore establish three points. First, commodity and exchange-rate relationships are economically relevant and empirically non-constant. Second, nonlinear and time-frequency methods often reveal features that are hidden in mean-based specifications. Third, there is already a domestic precedent for copula-based tail analysis. What remains less developed in this literature is the interaction between daily market-calendar construction, absolute copula adequacy, marginal misspecification and corrected inference across many tail comparisons. That narrower gap motivates the design adopted in this study.

\subsection{International commodity-currency dependence and tail risk}

The international commodity-currency literature also shows substantial heterogeneity. \citet{reboredo2012} uses copulas to model oil-price and exchange-rate comovement and finds generally weak dependence in the currencies studied, although dependence strengthened after the global financial crisis. \citet{ignatieva2017}, examining major commodity exporters at daily frequency, find positive dependence between commodity-price indices and commodity currencies and report a pronounced increase in time-varying tail dependence after the global financial crisis. These results are not contradictory so much as evidence that dependence is sensitive to country, commodity composition, sample period and model specification.

Later studies place greater emphasis on extreme risk rather than average comovement. \citet{bondatti2022} analyse downside commodity tail risk in exchange rates and show that currency exposure is heterogeneous across commodities, including oil and gold. \citet{ning2025} use a dependence-switching copula framework to study oil prices and exchange rates in both oil-exporting and oil-importing countries and report that dependence and downside/upside spillovers vary across regimes. Such results caution against treating a commodity currency as having one stable dependence coefficient that applies equally in calm and stressed markets.

Recent evidence extends that argument. \citet{go2024}, using Malaysian crude palm oil and foreign-exchange markets, document changes in return and volatility spillovers across crisis and post-crisis periods. \citet{dodd2026} show that the usually positive contemporaneous relationship between commodity and currency returns in major commodity-exporting countries can disappear when political risk rises. Their result is especially useful for the present paper because it shows that even a well-established commodity-currency relation can be state-dependent. The implication is not that Ghana must follow the same pattern, but that robustness to market conditions should be tested rather than assumed.

This literature also helps define what the present paper does not claim. It does not estimate causal transmission from commodity prices to the Cedi, and it does not attempt to identify a universal commodity-currency pricing mechanism. Instead, it asks whether observed dependence and tail-concentration conclusions are stable across two defensible calendar constructions and whether any of eight static bivariate copula families provides an adequate representation of those relationships. This narrower objective is appropriate given the heterogeneity documented in the broader literature.

\subsection{Crisis periods and changes in extreme dependence}

A large post-2020 literature reports stronger cross-market dependence during periods of financial distress. In commodity markets, \citet{qiao2023} use a copula-based network approach and document higher lower- and upper-tail contagiousness after the onset of COVID-19. Chen et al. (2023) likewise find that dependence among commodity futures changes substantially during major distress episodes when flexible dynamic copula models are used. Related quantile-connectedness work also reports stronger transmission in the tails during the pandemic.

These findings provide a strong reason to test stress-period changes, but they do not establish that every commodity pair or every currency relationship must show stronger tail dependence in a crisis. Samples differ in market coverage, frequency, stress-window definition, tail estimand and model structure. They also differ in how many hypotheses are evaluated and whether multiplicity is controlled. The present study therefore treats COVID-19 and 2024 as prespecified stress windows and asks whether finite-threshold lower- and upper-tail concentration differs from the calm baseline under the same inferential design in both calendar panels.

This distinction is important because nominal significance can accumulate quickly in a large set of pair, threshold and tail comparisons. \citet{bormann2020} provide a close methodological precedent in showing that structural differences in financial tail dependence can be studied with bootstrap procedures while accounting for multiple comparisons. More recent work develops formal two-sample tests for equality of tail copulas (Can, Einmahl and Laeven, 2024) and tests of tail equivalence (Koike, Kato and Yoshiba, 2025). These procedures address broader inferential objects than the finite-threshold concentration differences used here, but they reinforce a common principle: the statistical claim must match the tail quantity actually tested.

Accordingly, the present stress analysis is intentionally modest in interpretation. A failure to reject a stress-calm difference is not evidence that the two regimes are identical, and an unadjusted p-value below 0.05 is not treated as a robust discovery when the family contains many simultaneous tests. The relevant question is whether a result survives the prespecified multiplicity correction and whether that corrected conclusion agrees across both calendar constructions.

\subsection{Nonsynchronous trading and calendar construction}

The market-calendar problem has a long history in empirical finance. \citet{scholes1977} show that nonsynchronous trading can introduce serious errors-in-variables problems into return-based models. \citet{lo1990} formalize how infrequent or nonsynchronous trading affects variances, autocorrelations and cross-autocorrelations. These early studies establish that observed return dependence can reflect the observation process as well as the underlying economic relationship.

The problem becomes particularly visible when assets trade in different markets or at different times. \citet{martens2001} show that close-to-close international stock returns can understate correlation when trading hours differ and that estimated covariance and correlation dynamics change after synchronization. \citet{schotman2006}, studying Central European emerging markets, similarly find that time adjustment improves market-integration estimates and that aggregation to weekly frequency does not eliminate the time-matching problem. More recently, \citet{dimitriou2025} show that synchronization choices materially affect volatility-spillover inference in G-7 equity markets during crisis periods.

These studies do not prescribe one universally correct calendar treatment for the present data. Their contribution is more fundamental: calendar alignment is part of model construction and can change measured dependence. That is why this study uses two co-equal panels rather than designating one as the true dataset and the other as a robustness check. The business-day panel preserves the documented business-day construction with limited forward filling, while the common-observation panel computes returns only between consecutive dates on which all five series have actual observations. The resulting comparison asks which conclusions survive both plausible representations of the observation process.

The common-observation panel should not be interpreted as perfect intraday synchronization. Its return intervals range across consecutive common observation dates and can span more than one calendar day. The business-day panel, by contrast, preserves nominal business-day frequency but can generate zero returns when prices are carried forward. These are different compromises. Treating the two constructions symmetrically makes the consequences of those compromises visible rather than embedding them silently in a single preprocessing choice.

\subsection{Copula adequacy, marginal specification and tail inference}

Copulas separate marginal behavior from cross-variable dependence, which makes them attractive for financial returns with non-Gaussian margins. That separation does not make the marginal stage irrelevant. \citet{fantazzini2009} shows through simulation that misspecified margins can materially bias estimated copula parameters and risk measures. This is directly relevant to the Cedi in the present study because its return series contains many exact zeros and the selected continuous marginal model fails the same probability-integral-transform adequacy diagnostics passed by the commodity margins. The resulting Cedi-related copula and extreme-value quantities are therefore interpreted conditionally on an inadequate marginal reference model rather than as model-free structural facts.

The treatment of ties deserves similar care. \citet{kojadinovic2017} shows that many standard copula procedures are derived under the assumption of continuous margins without ties and can be biased when ties arise through rounding or limited measurement precision. The present project therefore audited ties explicitly. Panel B contains no exact PIT ties. Panel A contains one three-way Cedi PIT tie created by the numerical clipping of three extreme observations at \(10^{-6}\). Recovering their underlying rank order changes only two pseudo-observations, leaves membership in the prespecified \(q=.025,.05,.10\) lower-tail sets unchanged, and produces only negligible changes in the observed Cedi-pair goodness-of-fit statistics. The clipping issue is therefore treated as a closed numerical audit rather than evidence that all Cedi copula inference is invalid.

Goodness of fit also needs to be distinguished from relative model selection. \citet{genestremillard2008} establish conditions under which parametric bootstrap methods provide valid approximations for goodness-of-fit testing in semiparametric models. \citet{genest2009} review a broad class of copula goodness-of-fit procedures and emphasize bootstrap implementation under composite null hypotheses. \citet{dobric2007} study a goodness-of-fit procedure based on the Rosenblatt transformation and show that bootstrap treatment is important when margins are estimated rather than known. Consistent with this literature, the present paper reports a Rosenblatt-transform Cramér-von Mises statistic with parametric bootstrap. It does not equate the lowest AIC with model adequacy, and it does not describe the implemented statistic as an exact replication of every empirical-copula procedure considered by Genest, Rémillard and Beaudoin.

Tail measurement presents a separate problem. \citet{frahm2005} show that estimators of the asymptotic tail-dependence coefficient can behave differently in finite samples, while \citet{supper2020} show that static estimators can be severely biased when extreme dependence varies over time. These results support two wording choices in the present study. First, the empirical quantities calculated at \(q=.025,.05,.10\) are described as finite-threshold tail concentration measures rather than asymptotic tail-dependence coefficients. Second, full-sample static copula estimates are treated as summaries of the observed sample rather than as evidence that the underlying dependence process is time-invariant.

Finally, the large number of stress comparisons requires explicit multiplicity control. \citet{benjamini1995} introduced false-discovery-rate control as an alternative to family-wise error control, while Bonferroni adjustment provides a stricter family-wise criterion. The present analysis reports both. Each stress definition is treated as its own prespecified 60-test family, comprising ten pairs, three thresholds and two tails. This prevents isolated nominal p-values from being interpreted as robust stress effects when they do not survive the correction applied to the family from which they were selected.

\subsection{Positioning of the present study}

The literature reviewed above points to a specific rather than sweeping gap. Ghanaian studies already document asymmetric commodity-exchange-rate relationships, time-frequency connectedness and exchange-rate tail dependence. International studies already demonstrate that commodity-currency dependence can be nonlinear, dynamic and crisis-sensitive. Statistical research already provides formal tools for copula adequacy, tail comparison and multiple testing.

The present study connects these strands through a robustness question that is rarely made explicit in applied commodity-currency work: which conclusions survive alternative defensible treatments of nonsynchronous market calendars? It addresses that question using two co-equal panels, explicit marginal convergence and adequacy gates, eight static bivariate copula families with absolute goodness-of-fit testing, finite-threshold stress comparisons, and Bonferroni and Benjamini-Hochberg correction within prespecified test families. It also separates calendar-robust findings from calendar-sensitive and panel-specific diagnostics. The contribution is therefore not a claim that a new dependence method has been invented or that copulas have not previously been used in Ghana. It is a design for distinguishing conclusions that are stable to consequential data-construction choices from those that are not.

\subsection{Working reference register for this section}

\begin{itemize}
\item Archer, C., Owusu Junior, P., Adam, A. M., Asafo-Adjei, E., \& Baffoe, S. (2022). *Asymmetric Dependence Between Exchange Rate and Commodity Prices in Ghana*. Annals of Financial Economics, 17(2). DOI: 10.1142/S2010495222500129.
\item Atipaga, U.-F., \& Mensah, M. (2025). *Global uncertainties and the Ghanaian cedi: time-frequency and directional information flow analysis*. Cogent Economics \& Finance, 13(1), 2598191. DOI: 10.1080/23322039.2025.2598191.
\item Barson, Z., Owusu Junior, P., \& Adam, A. M. (2023). *Comovement between commodity returns in Ghana: the role of exchange rates*. Journal of Economic Structures, 12, 17. DOI: 10.1186/s40008-023-00312-z.
\item Benjamini, Y., \& Hochberg, Y. (1995). *Controlling the False Discovery Rate: A Practical and Powerful Approach to Multiple Testing*. Journal of the Royal Statistical Society Series B, 57(1), 289--300. DOI: 10.1111/j.2517-6161.1995.tb02031.x.
\item Boateng, E., Asafo-Adjei, E., Addison, A., Quaicoe, S., Yusuf, M. A., Abeka, M. J., \& Adam, A. M. (2022). *Interconnectedness among commodities, the real sector of Ghana and external shocks*. Resources Policy, 75, 102511. DOI: 10.1016/j.resourpol.2021.102511.
\item Bondatti, M., \& Rillo, G. (2022). *Commodity tail-risk and exchange rates*. Finance Research Letters, 47, 102937. DOI: 10.1016/j.frl.2022.102937.
\item Bormann, C., \& Schienle, M. (2020). *Detecting Structural Differences in Tail Dependence of Financial Time Series*. Journal of Business \& Economic Statistics, 38(2), 380--392. DOI: 10.1080/07350015.2018.1506343.
\item Can, S. U., Einmahl, J. H. J., \& Laeven, R. J. A. (2024). *Two-Sample Testing for Tail Copulas with an Application to Equity Indices*. Journal of Business \& Economic Statistics, 42(1), 147--159. DOI: 10.1080/07350015.2023.2166050.
\item Dimitriou, D., \& Kenourgios, D. (2025). *The implications of non-synchronous trading in G-7 financial markets*. International Journal of Finance \& Economics, 30(1), 689--709. DOI: 10.1002/ijfe.2936.
\item Dobrić, J., \& Schmid, F. (2007). *A goodness of fit test for copulas based on Rosenblatt's transformation*. Computational Statistics \& Data Analysis, 51(9), 4633--4642. DOI: 10.1016/j.csda.2006.08.012.
\item Dodd, O., Fernandez-Perez, A., \& Sosvilla-Rivero, S. (2026). *Political risk and commodity currencies*. Journal of Commodity Markets, 42, 100562. DOI: 10.1016/j.jcomm.2026.100562.
\item Fantazzini, D. (2009). *The effects of misspecified marginals and copulas on computing the value at risk: A Monte Carlo study*. Computational Statistics \& Data Analysis, 53(6), 2168--2188. DOI: 10.1016/j.csda.2008.02.002.
\item Frahm, G., Junker, M., \& Schmidt, R. (2005). *Estimating the tail-dependence coefficient: Properties and pitfalls*. Insurance: Mathematics and Economics, 37(1), 80--100. DOI: 10.1016/j.insmatheco.2005.05.008.
\item Genest, C., \& Rémillard, B. (2008). *Validity of the parametric bootstrap for goodness-of-fit testing in semiparametric models*. Annales de l'Institut Henri Poincaré Probabilités et Statistiques, 44(6), 1096--1127. DOI: 10.1214/07-AIHP148.
\item Genest, C., Rémillard, B., \& Beaudoin, D. (2009). *Goodness-of-fit tests for copulas: A review and a power study*. Insurance: Mathematics and Economics, 44(2), 199--213. DOI: 10.1016/j.insmatheco.2007.10.005.
\item Go, Y.-H., \& Lau, W.-Y. (2024). *Terms of trade or market power? Further evidence from dynamic spillovers in return and volatility between Malaysian crude palm oil and foreign exchange markets*. North American Journal of Economics and Finance, 73, 102178. DOI: 10.1016/j.najef.2024.102178.
\item Ignatieva, K., \& Ponomareva, N. (2017). *Commodity currencies and commodity prices: modelling static and time-varying dependence*. Applied Economics, 49(15), 1491--1512. DOI: 10.1080/00036846.2016.1221038.
\item Kojadinovic, I. (2017). *Some copula inference procedures adapted to the presence of ties*. Computational Statistics \& Data Analysis, 112, 24--41. DOI: 10.1016/j.csda.2017.02.006.
\item Koike, T., Kato, S., \& Yoshiba, T. (2025). *Measuring and testing tail equivalence*. Journal of Multivariate Analysis, 209, 105460. DOI: 10.1016/j.jmva.2025.105460.
\item Lo, A. W., \& MacKinlay, A. C. (1990). *An econometric analysis of nonsynchronous trading*. Journal of Econometrics, 45(1--2), 181--211. DOI: 10.1016/0304-4076(90)90098-E.
\item Martens, M., \& Poon, S.-H. (2001). *Returns synchronization and daily correlation dynamics between international stock markets*. Journal of Banking \& Finance, 25(10), 1805--1827. DOI: 10.1016/S0378-4266(00)00159-X.
\item Mensah, P. O., \& Adam, A. M. (2020). *Copula-Based Assessment of Co-Movement and Tail Dependence Structure Among Major Trading Foreign Currencies in Ghana*. Risks, 8(2), 55. DOI: 10.3390/risks8020055.
\item Ning, C., \& Xu, D. (2025). *Extreme comovements and downside/upside risk spillovers between oil prices and exchange rates*. Macroeconomic Dynamics, 29, e53. DOI: 10.1017/S1365100524000543.
\item Qiao, T., \& Han, L. (2023). *COVID-19 and tail risk contagion across commodity futures markets*. Journal of Futures Markets, 43(2), 242--272. DOI: 10.1002/fut.22388.
\item Reboredo, J. C. (2012). *Modelling oil price and exchange rate co-movements*. Journal of Policy Modeling, 34(3), 419--440. DOI: 10.1016/j.jpolmod.2011.10.005.
\item Scholes, M., \& Williams, J. (1977). *Estimating betas from nonsynchronous data*. Journal of Financial Economics, 5(3), 309--327. DOI: 10.1016/0304-405X(77)90041-1.
\item Schotman, P. C., \& Zalewska, A. (2006). *Non-synchronous trading and testing for market integration in Central European emerging markets*. Journal of Empirical Finance, 13(4--5), 462--494. DOI: 10.1016/j.jempfin.2006.04.002.
\item Supper, H., Irresberger, F., \& Weiß, G. (2020). *A comparison of tail dependence estimators*. European Journal of Operational Research, 284(2), 728--742. DOI: 10.1016/j.ejor.2019.12.041.
\item Woode, J. K., Idun, A. A.-A., Kawor, S., Owusu Junior, P., \& Adam, A. M. (2025). *Dynamic Connectedness Between Commodities, Exchange Rates and Equity Markets of Commodity-Dependent Sub-Saharan Africa Countries*. SAGE Open, 15(3). DOI: 10.1177/21582440251347723.
\item Yeboah, S., Fumey, M., Winful, S. A., Otoo, I., \& Owusu, P. (2025). *Asymmetric dependence between commodity prices and selected macroeconomic variables in Ghana*. Scientific African, 28, e02739. DOI: 10.1016/j.sciaf.2025.e02739.
\item Zankawah, M. M., \& Stewart, C. (2020). *Measuring the volatility spill-over effects of crude oil prices on the exchange rate and stock market in Ghana*. Journal of International Trade \& Economic Development, 29(4), 420--439. DOI: 10.1080/09638199.2019.1692895.

\end{itemize}
\textbf{Note:} Bibliographic metadata will be normalized against Crossref and the target journal style before conversion to the final \texttt{references.bib}.
